% !TeX root = ../report.tex

\makeatletter

\newcommand{\course}[1]{\def\@course{#1}}
\newcommand{\group}[1]{\def\@group{#1}}
\newcommand{\theme}[1]{\def\@theme{#1}}
\newcommand{\studentid}[1]{\def\@studentid{#1}}
\newcommand{\experimentpartner}[1]{\def\@experimentpartner{#1}}
\newcommand{\staff}[1]{\def\@staff{#1}}
\newcommand{\experimentdate}[1]{\def\@experimentdate{#1}}
\newcommand{\place}[1]{\def\@place{#1}}
\newcommand{\experimenttime}[1]{\def\@experimenttime{#1}}
\newcommand{\reportfirstdue}[1]{\def\@reportfirstdue{#1}}

\providecommand{\@course}{}
\providecommand{\@group}{}
\providecommand{\@theme}{}
\providecommand{\@studentid}{}
\providecommand{\@experimentpartner}{}
\providecommand{\@staff}{}
\providecommand{\@experimentdate}{}
\providecommand{\@place}{}
\providecommand{\@experimenttime}{}
\providecommand{\@reportfirstdue}{}

\renewcommand{\maketitle}{%
  \begin{center}
    {\Large \bfseries レポートの基本情報\par}
    \vspace{1em}
    \begin{tabularx}{0.95\linewidth}{@{}>{\raggedright\arraybackslash}p{0.25\linewidth}
                                    >{\centering\arraybackslash}m{1em}
                                    >{\raggedright\arraybackslash}X@{}}
      \hline
      実験テーマ名 & ： & \@title \\
      レポート執筆者 & ： & \@author \\
      共同実験者 & ： & \@experimentpartner \\
      担当教員やTA名 & ： & \@staff \\
      実験実施日 & ： & \@experimentdate \\
      実施場所 & ： & \@place \\
      時間 & ： & \@experimenttime \\
      初回提出日 & ： & \@reportfirstdue \\
      \hline
    \end{tabularx}
  \end{center}
}

\pagestyle{fancy}
\fancyhf{}
\fancyhead[C]{%
  \small
  \begingroup
  \renewcommand{\arraystretch}{0.95}%
  \begin{tabularx}{\textwidth}{@{}
      >{\raggedright\arraybackslash}m{0.26\textwidth}
      >{\centering\arraybackslash}m{0.10\textwidth}
      >{\centering\arraybackslash}m{0.19\textwidth}
      >{\centering\arraybackslash}m{0.08\textwidth}
      >{\raggedleft\arraybackslash}X@{}}
    \@course & テーマ\ \@theme & 学籍番号：\@studentid & \@group 班 & 作成者：\@author
  \end{tabularx}%
  \endgroup
}
\fancyfoot[C]{\thepage}
\makeatother
